\chapter{Discussion}
\label{ch:discussion}

\section{Interpretation of Results}

\subsection{OCR Ensemble Effectiveness}

The ensemble's character accuracy of~89.1\% is notably higher than either individual
engine (78.1\% and 83.4\%), a relative improvement of~14.1\% over Tesseract and~6.8\%
over EasyOCR. This is consistent with the ensemble literature: engines trained on different
data distributions with different architectures tend to fail on complementary subsets of the
input~\cite{ratzlaff2003method}. In the IKEA label domain, Tesseract's strengths lie in
high-contrast numeric fields (article numbers, barcodes' human-readable supplements) while
EasyOCR is stronger on smaller serif text and multi-line blocks.

The 6-point improvement comes at a cost of~260~ms additional latency. For the intended
use case—pre-press validation, not real-time inspection—this is an acceptable trade-off.

\subsection{Label-Type Classifier Accuracy}

An 86.1\% top-1 accuracy without any trained model is noteworthy. The scoring heuristic
(\cref{ch:system}, Eq.~\ref{eq:type_score}) is simple by design—it counts field presences
and adds barcode bonuses—yet it correctly identifies the label type for 31 of 36 test labels.
The five misclassifications all fall within the same label family, suggesting the heuristic
lacks the discriminating signal needed to distinguish intra-family variants.

A natural extension would be to add physical-dimension inference: if the aspect ratio of
the label image is consistent with a known label size (e.g.~70~mm~$\times$~35~mm for 1R70
vs.~90~mm~$\times$~35~mm for 1R90), this can disambiguate many intra-family cases. This
requires a known physical reference (e.g.\ a calibration target or a known barcode module
size), which is available in production but not always in the test dataset.

\subsection{Field Extraction Challenges}

The three lowest-scoring fields—\texttt{plant\_identifier} (F1~=~0.794),
\texttt{dimensions\_metric} (F1~=~0.804), and \texttt{gross\_weight} (F1~=~0.820)—share
a common characteristic: their values are formatted inconsistently across label types. The
plant identifier pattern \texttt{PI-XXX-YYY} is sometimes separated by en-dashes rather
than hyphens; dimension strings use both ``cm'' and ``mm'' units with varying decimal separators
depending on the target market. Future work should expand the regex patterns to cover a
wider range of formatting variants observed in real production data.

\subsection{Validation Precision vs.\ Recall Trade-off}

A precision of~88.3\% and recall of~84.7\% at the default threshold means that:
\begin{itemize}
  \item 2 of 16 non-compliant labels in the test set were missed (false negatives). Both
  involved violations not covered by any rule in the current schema.
  \item 2 of 20 compliant labels were incorrectly flagged (false positives). Both involved
  labels where an optional field was absent and the rule schema had incorrectly marked it
  as required.
\end{itemize}

The false positive cases highlight an important limitation: the schema accuracy is bounded
by the accuracy of the human expert who authored the JSON rule files. Two rules in the
\texttt{1C70-IDADM} schema were found post-hoc to list fields as required that the IKEA
specification actually marks as conditional on other fields being present. Fixing these two
rules would raise accuracy from~88.9\% to~94.4\% on the current test set.

\section{Limitations}

\subsection{Rule Corpus Incompleteness}

The 34 schemas cover~342 required-field declarations and~107 rules, but the design
specification contains additional structural requirements that are difficult to encode as
simple field-level rules:

\begin{itemize}
  \item \textbf{Zone ordering}: the canonical left-to-right zone order must be preserved.
  The layout checker stub (\texttt{layout\_checker.py}) is designed to implement this but
  was not completed within the project timeline.
  \item \textbf{Visual hierarchy}: product name font must be larger than article number
  font. This requires font-size inference from pixel heights, which is fragile without
  a physical resolution calibration.
  \item \textbf{Conditional presence}: some fields are required only if other fields are
  present (e.g.\ the YYWW stamp is required only if the date zone is activated). The
  current schema format does not support conditional rules.
\end{itemize}

\subsection{OCR Limitations on Rotated Text}

Several IKEA label types (particularly the 1R35 and 1R10 families) contain text zones
rotated~90° clockwise—typically the compliance block and copy-count strip. The current
preprocessing applies a global image orientation but does not detect and deskew individual
zones independently. EasyOCR handles moderate rotation through its CRAFT detector, but
Tesseract with \texttt{psm 11} does not. A future improvement is to detect rotated zones
and apply independent rotation correction before the OCR pass.

\subsection{DataMatrix Decode Rate at Small Magnification}

The 88.7\% overall DataMatrix decode rate masks a significant gap between label types:
for labels with DataMatrix at~0.4~mm x-dim (the smallest in the IKEA catalogue), the
decode rate at 300~DPI drops to~74\%. For labels with~0.75~mm x-dim, the decode rate
is~97\%. Improving small-DataMatrix reliability requires either higher rendering DPI (400+)
or a super-resolution preprocessing step.

\subsection{Dataset Size}

The 238-label dataset is small relative to the diversity of the label catalogue. Six label
types have no representation, and the represented types average fewer than nine samples each.
This makes it difficult to draw strong conclusions about performance on under-represented
types. Expanding the dataset—particularly with real production failures collected by IKEA's
\ac{QA} team—is the most impactful next step for improving reliability.

\section{Ethical Considerations}

The system is designed to support, not replace, human QA reviewers. All validation
decisions are presented with a full rule-by-rule trace, ensuring that a human can audit
any decision the system makes. The system does not store or transmit label content to
external services; all processing is on-premise.

One concern is automation bias: a QA reviewer who trusts the system too much may fail
to catch violations that the system misses. This risk is mitigated by the explicit listing
of uncovered rule types (layout and completeness) in the UI and by displaying the
\texttt{compliance\_score} as a single number rather than a binary pass/fail verdict,
encouraging reviewers to exercise judgement when the score is in the 0.55--0.85 range.

\section{Industry Applicability}

The approach taken in this project—encoding domain expertise as JSON schema rules and
applying them to extracted structured data—is deliberately technology-conservative. The
pipeline does not require a GPU, does not require training data to add new label types
(only a new JSON schema), and does not depend on any cloud service. These properties make
it practical for IKEA of Sweden's Älmhult facility, which requires on-premise processing
for data security reasons.

The estimated processing rate of~1~label per~1.4~seconds (ensemble mode) is sufficient
for the target use case of pre-press batch validation. A typical batch of~200~labels before
a product launch would complete in under~5 minutes, compared to an estimated~26 hours of
manual review at the current rate of 50--70 labels per working day.
